% (c) Nikita Lisitsa, lisyarus@gmail.com, 2026

\documentclass[10pt]{beamer}

\usepackage[T2A]{fontenc}
\usepackage[russian]{babel}
\usepackage{minted}

\usepackage[most]{tcolorbox}
\tcbuselibrary{minted}

\usepackage{graphicx}
\graphicspath{ {./images/} }

\usepackage{adjustbox}

\usepackage{color}
\usepackage{soul}
\usepackage{multicol}
\usepackage{multirow}
\usepackage{hyperref}
\usepackage{tikz}
\usetikzlibrary{arrows.meta,positioning,calc}

\usetheme{metropolis}

\definecolor{red}{rgb}{1,0,0}
\definecolor{green}{rgb}{0,0.5,0}
\definecolor{codebg}{RGB}{29,35,49}
\setminted{fontsize=\footnotesize}

\makeatletter
\newcommand{\slideimage}[1]{
  \begin{figure}
    \begin{adjustbox}{width=\textwidth, totalheight=\textheight-2\baselineskip-2\baselineskip,keepaspectratio}
      \includegraphics{#1}
    \end{adjustbox}
  \end{figure}
}
\makeatother

\title{Язык программирования C}
\subtitle{Лекция 7: inline, variadic functions, статические и динамические библиотеки}
\author{Лисица Никита}
\date{2026}

\setbeamertemplate{footline}[frame number]

\usemintedstyle{tango}

\begin{document}

\frame{\titlepage}

% \begin{frame}[fragile]
% \frametitle{Инлайнинг}
% \usemintedstyle{lightbulb}
% \begin{tcblisting}{listing engine=minted, minted language=cpp, minted options={fontsize=\scriptsize}, listing only, colback=codebg, colframe=codebg, rounded corners}
% int* get_array() {
%   int array[100];
%   for(int i = 0; i < 100; ++i) {
%     array[i] = i;
%   }
%   return array;
% }

% int main() {
%   int *array = get_array();
%   printf("%d", array[0]);
% }
% \end{tcblisting}
% \usemintedstyle{tango}
% \begin{itemize}
% \item Что в этом коде не так? \pause \mintinline{cpp}|array| -- локальная переменная: она живёт на стеке, и будет ``удалена'' после выхода из функции \mintinline{cpp}|get_array()|
% \end{itemize}
% \end{frame}

\begin{frame}[fragile]
\frametitle{Инлайнинг}
\begin{itemize}
\item В коде часто встречаются маленькие вспомогательные функции:
\end{itemize}
\usemintedstyle{lightbulb}
\begin{multicols}{2}
\begin{tcblisting}{listing engine=minted, minted language=cpp, minted options={fontsize=\scriptsize}, listing only, colback=codebg, colframe=codebg, rounded corners}
int max(int x, int y) {
  return (x < y) ? y : x;
}
\end{tcblisting}
\columnbreak
\begin{tcblisting}{listing engine=minted, minted language=cpp, minted options={fontsize=\scriptsize}, listing only, colback=codebg, colframe=codebg, rounded corners}
float sqr(float x) {
  return x * x;
}
\end{tcblisting}
\end{multicols}
\pause
\begin{itemize}
\item При вызове такой функции больше времени тратится на сам вызов, чем на исполнение кода функции
\end{itemize}
\end{frame}

\begin{frame}[fragile]
\frametitle{Инлайнинг}
\begin{itemize}
\item \textit{Инлайнинг (inlining)} -- оптимизация компилятора, при которой код функции подставляется целиком в место её вызова (in-line)
\end{itemize}
\pause
\usemintedstyle{lightbulb}
\begin{multicols}{2}
\begin{tcblisting}{listing engine=minted, minted language=cpp, minted options={fontsize=\scriptsize}, listing only, colback=codebg, colframe=codebg, rounded corners}
// Было
int a = 10, b = 20;
int x = max(a, b);
\end{tcblisting}
\columnbreak
\begin{tcblisting}{listing engine=minted, minted language=cpp, minted options={fontsize=\scriptsize}, listing only, colback=codebg, colframe=codebg, rounded corners}
// Стало
int a = 10, b = 20;
int x = (a < b) ? b : a;
\end{tcblisting}
\end{multicols}
\usemintedstyle{tango}
\pause
\begin{itemize}
\item Компилятор \textbf{сам решает}, когда делать эту оптимизацию (зависит от включенного уровня оптимизации)
\pause
\item Можно форсировать специфичными для компилятора механизмами: \mintinline{cpp}|__attribute__((always_inline))| в GCC, \mintinline{cpp}|__forceinline| в MSVC
\end{itemize}
\end{frame}

\begin{frame}[fragile]
\frametitle{Инлайнинг}
\begin{itemize}
\item Мы привыкли объявлять функции в заголовочном файле и описывать их код в \mintinline{cpp}|.c|-файле
\end{itemize}
\usemintedstyle{lightbulb}
\begin{multicols}{2}
\begin{tcblisting}{listing engine=minted, minted language=cpp, minted options={fontsize=\scriptsize, escapeinside=||}, listing only, colback=codebg, colframe=codebg, rounded corners}
|{\phantom{Xg}}|
// max.h
int max(int a, int b);
|{\phantom{Xg}}|
\end{tcblisting}
\columnbreak
\begin{tcblisting}{listing engine=minted, minted language=cpp, minted options={fontsize=\scriptsize}, listing only, colback=codebg, colframe=codebg, rounded corners}
// max.c
int max(int x, int y) {
  return (x < y) ? y : x;
}
\end{tcblisting}
\end{multicols}
\usemintedstyle{tango}
\pause
\begin{itemize}
\item Сможет ли компилятор в этом случае заинлайнить вызов функции вне файла \mintinline{cpp}|max.c|? \pause Нет -- для инлайнинга нужен доступ к коду, а не только к объявлению
\end{itemize}
\end{frame}

\begin{frame}[fragile]
\frametitle{Инлайнинг}
\begin{itemize}
\item Решение: поместим код функции прямо в заголовочный файл:
\end{itemize}
\usemintedstyle{lightbulb}
\begin{tcblisting}{listing engine=minted, minted language=cpp, minted options={fontsize=\scriptsize}, listing only, colback=codebg, colframe=codebg, rounded corners}
// max.h
int max(int a, int b) {
  return (x < y) ? y : x;
}
\end{tcblisting}
\usemintedstyle{tango}
\pause
\begin{itemize}
\item Есть ли с этим проблемы? \pause Да: multiple definition, если \mintinline{cpp}|max.h| заинклюдят хотя бы два других файла
\end{itemize}
\usemintedstyle{lightbulb}
\begin{multicols}{2}
\begin{tcblisting}{listing engine=minted, minted language=cpp, minted options={fontsize=\scriptsize}, listing only, colback=codebg, colframe=codebg, rounded corners}
// a.c
#include <max.h>

...
\end{tcblisting}
\columnbreak
\begin{tcblisting}{listing engine=minted, minted language=cpp, minted options={fontsize=\scriptsize}, listing only, colback=codebg, colframe=codebg, rounded corners}
// b.c
#include <max.h>

...
\end{tcblisting}
\end{multicols}
\end{frame}

\begin{frame}[fragile]
\frametitle{Атрибут inline}
\begin{itemize}
\item Чтобы обойти эту проблему, у функций есть особый атрибут \mintinline{cpp}|inline|:
\end{itemize}
\usemintedstyle{lightbulb}
\begin{tcblisting}{listing engine=minted, minted language=cpp, minted options={fontsize=\scriptsize}, listing only, colback=codebg, colframe=codebg, rounded corners}
// max.h
inline int max(int a, int b) {
  return (x < y) ? y : x;
}
\end{tcblisting}
\usemintedstyle{tango}
\pause
\begin{itemize}
\item \mintinline{cpp}|inline| -- указание \textbf{линковщику}, что multiple definition этой конкретной фукнции -- \textbf{не ошибка}, и он может взять любую
\pause
\item Если у вас есть несколько \textit{разных} \mintinline{cpp}|inline|-функций с одинаковым названием, это undefined behavior
\end{itemize}
\end{frame}

\begin{frame}[fragile]
\frametitle{Атрибут inline}
\begin{itemize}
\item \(\Longrightarrow\) Если вы хотите, чтобы функция инлайнилась, нужно
\begin{enumerate}
\item Положить её реализацию в заголовочный файл, и
\item Пометить её атрибутом \mintinline{cpp}|inline|, чтобы не ругался линковщик
\end{enumerate}
\pause
\item \textbf{N.B.:} На решение компилятора делать инлайнинг атрибут \mintinline{cpp}|inline| \textbf{не влияет}, он сам разберётся :)
\pause
\item \mintinline{cpp}|inline| -- флаг именно для линковщика
\end{itemize}
\end{frame}

\begin{frame}[fragile]
\frametitle{Variadic functions}
\begin{itemize}
\item Мы видели функции, принимающие произвольное количество аргументов: \mintinline{cpp}|printf|, \mintinline{cpp}|scanf|, и т.п. Как это работает?
\pause
\item В объявлении таких функций после списка аргументов стоит \textit{многочие}, соответствующее \textit{variadic} аргументам: \mintinline{cpp}|int printf(const char *format, ...);|
\pause
\item Концептуально: variadic аргументы функции кладутся на стек, и ожидается, что функция сама знает их количество и типы (например, \mintinline{cpp}|printf| знает это из строки формата)
\pause
\item На практике аргументы передаются частично через регистры и частично через стек (см. ABI и calling conventions конкретной платформы)
\pause
\item \textbf{N.B.:} В C++ мы встретим \textit{variadic templates} -- это похожий по общее идее, но совершенно другой по реализации механизм :)
\end{itemize}
\end{frame}

\begin{frame}[fragile]
\frametitle{Variadic functions}
\begin{itemize}
\item Для реализации variadic функций есть специальные типы и функции (на самом деле -- макросы):
\pause
\begin{itemize}
\item \mintinline{cpp}|va_start()|, \mintinline{cpp}|va_end()| -- начать и закончить обработку аргументов
\pause
\item \mintinline{cpp}|va_arg()| -- получить текущий аргумент и перейти к следующему
\pause
\item \mintinline{cpp}|va_list| -- объект, хранящий состояние прохода по аргументам
\end{itemize}
\end{itemize}
\end{frame}

\begin{frame}[fragile]
\frametitle{Variadic functions}
\begin{itemize}
\item Пример variadic функции, реализующей сумму произвольного набора чисел (первый аргумент -- количество):
\end{itemize}
\usemintedstyle{lightbulb}
\begin{tcblisting}{listing engine=minted, minted language=cpp, minted options={fontsize=\scriptsize}, listing only, colback=codebg, colframe=codebg, rounded corners}
int sum(int count, ...) {
  int result = 0;
  va_list args;
  va_start(args, count);
  while (count > 0) {
    result += va_arg(args, int);
  }
  va_end(args);
  return result;
}
\end{tcblisting}
\end{frame}

\begin{frame}[fragile]
\frametitle{Variadic functions}
\begin{itemize}
\item Пример variadic функции, реализующей примитивный printf:
\end{itemize}
\usemintedstyle{lightbulb}
\begin{tcblisting}{listing engine=minted, minted language=cpp, minted options={fontsize=\scriptsize}, listing only, colback=codebg, colframe=codebg, rounded corners}
void printf(const char *format, ...) {
  va_list args;
  va_start(args, format);
  while (*format != '\0') {
    if (*format == '%') {
      ++format;
      if (*format == 'd') {
        int value = va_arg(args, int);
        // Печатаем int
      } else if (*format == 's') {
        const char * value = va_arg(args, const char *);
        // Печатаем строку
      }
    } else {
      char c = *format;
      // Печатаем char
    }
    ++format;
  }
  va_end(args);
}
\end{tcblisting}
\end{frame}

\begin{frame}[fragile]
\frametitle{Variadic functions: проблемы}
\begin{itemize}
\item Трудно обнаружить ошибки несоответствия типов/количества аргументов:
\end{itemize}
\usemintedstyle{lightbulb}
\begin{tcblisting}{listing engine=minted, minted language=cpp, minted options={fontsize=\scriptsize}, listing only, colback=codebg, colframe=codebg, rounded corners}
// Все эти вызовы - undefined behavior:
printf("%d %d", 10);
printf("%d", 11, 12);
printf("%s", 3.1415f);
\end{tcblisting}
\usemintedstyle{tango}
\pause
\begin{itemize}
\item Современные компиляторы и IDE знают про библиотечные variadic функции (\mintinline{cpp}|printf| и т.п.) и выдают предупреждения при несоответствии типов/количества аргументов
\pause
\item Механизм \textit{variadic templates} в C++ (про который мы узнаем позже) лишён этого недостатка
\end{itemize}
\end{frame}

\begin{frame}[fragile]
\frametitle{Виды библиотек}
\begin{itemize}
\item \textit{Библиотека (library)} -- набор кода, предназначенного для переиспользования в разных программах/проектах
\pause
\item Они различаются по тому, как мы их подключаем к нашему коду:
\pause
\begin{itemize}
\item \textit{Header-only} библиотеки
\pause
\item \textit{Статический (static)} библиотеки
\pause
\item \textit{Динамические (dynamic)} библиотеки
\end{itemize}
\end{itemize}
\end{frame}

\begin{frame}[fragile]
\frametitle{Header-only библиотеки}
\begin{itemize}
\item \textit{Header-only} библиотеки состоят только из заголовочных файлов
\pause
\item Для их подключения достаточно сказать компилятору правильный путь до заголовочных файлов (\mintinline{text}|-I| в GCC, и т.п.)
\pause
\item Проще всего для добавления в существующую систему сборки
\pause
\item Если библиотека используется несколькими программами, её приходится компилировать каждый раз заново
\end{itemize}
\end{frame}

\begin{frame}[fragile]
\frametitle{Header-only библиотеки: пример}
\usemintedstyle{lightbulb}
\begin{tcblisting}{listing engine=minted, minted language=cpp, minted options={fontsize=\scriptsize}, listing only, colback=codebg, colframe=codebg, rounded corners}
// vector.h

typedef struct {
  float x, y, z;
} vec3;

vec3 add(vec3 a, vec3 b) {
  return (vec3){a.x + b.x, a.y + b.y, a.z + b.z};
}

vec3 muls(vec3 a, float b) {
  return (vec3){a.x * b, a.y * b, a.z * b};
}
\end{tcblisting}
\end{frame}

\begin{frame}[fragile]
\frametitle{Header-only библиотеки: пример}
\usemintedstyle{lightbulb}
\begin{tcblisting}{listing engine=minted, minted language=cpp, minted options={fontsize=\scriptsize}, listing only, colback=codebg, colframe=codebg, rounded corners}
// main.c
#include <vector.h>

int main() {
  vec3 v = {1.f, 2.f, 3.f};
  vec3 u = add(v, muls(v, 0.5f));
  return 0;
}
\end{tcblisting}
\begin{tcblisting}{listing engine=minted, minted language=bash, minted options={fontsize=\scriptsize}, listing only, colback=codebg, colframe=codebg, colupper=white, rounded corners}
gcc main.c -o main -I path/to/vector/includes
\end{tcblisting}
\end{frame}

\begin{frame}[fragile]
\frametitle{Статические библиотеки}
\begin{itemize}
\item \textit{Статические (static)} библиотеки состоят из заголовочных и \mintinline{text}|.c| файлов
\pause
\item Заголовочные файлы описывают \textit{публичный интерфейс} библиотеки (функции, через которые мы пользуемся этой библиотекой)
\pause
\item \mintinline{text}|.c| файлы компилируются в бинарный файл библиотеки (\mintinline{text}|.a| в Linux, \mintinline{text}|.lib| в Windows)
\pause
\item При \textbf{линковке нашей программы} статическая библиотека \textit{склеивается} с нашей программой
\pause
\item При обновлении библиотеки придётся перекомпилировать программу
\end{itemize}
\end{frame}

\begin{frame}[fragile]
\frametitle{Статические библиотеки: пример}
\usemintedstyle{lightbulb}
\begin{tcblisting}{listing engine=minted, minted language=cpp, minted options={fontsize=\scriptsize}, listing only, colback=codebg, colframe=codebg, rounded corners}
// vector.h

typedef struct {
  float x, y, z;
} vec3;

vec3 add(vec3 a, vec3 b);
vec3 muls(vec3 a, float b);
\end{tcblisting}
\usemintedstyle{lightbulb}
\begin{tcblisting}{listing engine=minted, minted language=cpp, minted options={fontsize=\scriptsize}, listing only, colback=codebg, colframe=codebg, rounded corners}
// vector.c
#include <vector.h>

vec3 add(vec3 a, vec3 b) {
  return (vec3){a.x + b.x, a.y + b.y, a.z + b.z};
}

vec3 muls(vec3 a, float b) {
  return (vec3){a.x * b, a.y * b, a.z * b};
}
\end{tcblisting}
\end{frame}

\begin{frame}[fragile]
\frametitle{Статические библиотеки: пример}
\usemintedstyle{lightbulb}
\begin{tcblisting}{listing engine=minted, minted language=bash, minted options={fontsize=\scriptsize}, listing only, colback=codebg, colframe=codebg, colupper=white, rounded corners}
# Компиляция библиотеки
gcc -c vector.c -o vector.o \
  -I path/to/vector/includes

# "Архивация библиотеки"
ar rcs libvector.a vector.o

# Компиляция программы
gcc main.c -o main -lvector \
  -I path/to/vector/includes \
  -L path/to/vector/lib
\end{tcblisting}
\end{frame}

\begin{frame}[fragile]
\frametitle{Динамические библиотеки}
\begin{itemize}
\item \textit{Динамические (static)} библиотеки компилируются независимо от нашей программы и лежат в отдельных файлах (\mintinline{text}|.so| в Linux, \mintinline{text}|.dll| в Windows, \mintinline{text}|.dylib| в macOS)
\pause
\item Могут устанавливаться вместе с программой, а могут поставляться системой (и устанавливаться напр. через менеджеры пакетов)
\pause
\item Наша программа \textbf{при запуске} загружает нужные динамические библиотеки в своё адресное пространство \(\Longrightarrow\) программе нужно как-то узнать, где они лежат (в Linux за это отвечает механизм RPATH)
\end{itemize}
\end{frame}

\begin{frame}[fragile]
\frametitle{Динамические библиотеки}
\begin{itemize}
\item Вызов функции из динамической библиотеки чуть дороже, чем вызов обычной функции (см. Procedure Linkage Table и Global Offset Table в ELF)
\pause
\item При изменении кода (но не публичного интерфейса!) динамической библиотеки достаточно просто перекомпилировать саму библиотеку -- использующие её программы автоматически загрузят новую версию при запуске
\pause
\item Почти всегда системные библиотеки, стандартная библиотека C, и т.п. -- динамические
\end{itemize}
\end{frame}

\begin{frame}[fragile]
\frametitle{Динамические библиотеки: пример}
\usemintedstyle{lightbulb}
\begin{tcblisting}{listing engine=minted, minted language=bash, minted options={fontsize=\scriptsize}, listing only, colback=codebg, colframe=codebg, colupper=white, rounded corners}
# Компиляция библиотеки
gcc -fPIC -c vector.c -o vector.o \
  -I path/to/vector/includes

# Линковка библиотеки
gcc -shared vector.o -o libvector.so

# Компиляция программы
gcc main.c -o main -lvector \
  -I path/to/vector/includes \
  -L path/to/vector/lib \
  -Wl,-rpath,path/to/vector/lib
\end{tcblisting}
\end{frame}

\end{document}
